\documentclass{article}

\IfFileExists{neurips_2027.sty}{\usepackage{neurips_2027}}{\usepackage[margin=1in]{geometry}}
\usepackage[utf8]{inputenc}
\usepackage[T1]{fontenc}
\usepackage{amsmath,amssymb,amsthm}
\usepackage{booktabs}
\usepackage{graphicx}
\usepackage{microtype}
\usepackage{natbib}
\usepackage{url}
\usepackage{xcolor}
\usepackage{hyperref}

\newtheorem{theorem}{Theorem}
\newtheorem{proposition}{Proposition}
\newcommand{\TRef}{\mathrm{ref}}
\newcommand{\Emu}{\mathrm{emu}}
\newcommand{\Hw}{\mathrm{hw}}
\newcommand{\Dis}{\mathsf{Dis}}
\newcommand{\TODO}[1]{\textcolor{red}{[TODO: #1]}}

\title{Transport Certificates for Spiking Neural Networks}
\author{Anonymous Authors}

\begin{document}
\maketitle

\begin{abstract}
Intermediate representations can move a trained spiking neural network (SNN)
between simulators and neuromorphic devices, but they do not answer the
deployment question: will the transported network preserve its decisions under
the target's discretization, reset, timing, numeric, and delivery semantics?
We formulate this question as prospective certification over a versioned family
of digital execution semantics. The proposed hierarchy combines exact semantic
mappings, sound finite-horizon per-input analysis, exact label-free statistical
bounds, and an emulator-to-hardware conformance term. The same certificate
directs label-free repair on a disjoint calibration split before an untouched
audit is certified. We describe realizations for SpiNNaker2 and a parameterized
Virtex-7 recurrent SNN pipeline, together with a preregistered evaluation on SHD,
DVS Gesture, and N-MNIST. Across five-seed software studies, disagreement ranks
semantic failures reliably and is tight on a feedforward negative control, but
its distribution-free accuracy bound is operationally loose on recurrent SHD
and convolutional/recurrent DVS Gesture. A sound interval analysis certifies
38.7\% of SHD inputs over a finite 16-member family, exactly matching enumerated
agreement after mutually exclusive semantics are checked separately; feedforward
N-MNIST reaches 98.6\%. However, adding only $\pm1\%$ continuous timestep and
threshold uncertainty makes the SHD interval certificate vacuous on every input,
even though 36.9\% agree on all tested box corners. Thus finite-family transport
is supported, but the intended bounded-family claim is not. Sequestered
replication and physical certificates remain unfinished.
\end{abstract}

\section{Introduction}

Neuromorphic intermediate representations make model structure portable, but
execution remains platform-specific. NIR intentionally describes continuous-time
primitives without fixing a digital discretization, and its cross-platform study
shows that differences in neuron definitions, integration, threshold timing,
quantization, and nondeterminism can change deployed dynamics
\citep{pedersen2024nir}. Recurrent connections can amplify a local transition
difference over time. A successful conversion therefore establishes that a graph
can be executed, not that its predictions have been preserved.

This distinction creates a practical decision problem. A deployer has trained
weights, an executable reference, a candidate backend, and unlabeled examples
from the intended distribution. Before paying for a full deployment---and before
revealing evaluation labels---the deployer wants to know whether the target's
absolute accuracy change is below a stated budget. If not, the deployer wants a
repair that changes only permitted neuron or numeric parameters and can itself be
re-certified.

Prediction disagreement is a natural observable, but it is not a new idea.
Disagreement has been analyzed for label-free performance estimation under
distribution shift \citep{lee2023dotl}, and disagreement discrepancy has been used
to derive error bounds \citep{rosenfeld2023disagreement}. Our setting uses a
simpler deterministic inequality: changing correctness on any labeled example
requires changing the prediction. The scientific question is consequently not
whether disagreement correlates with loss. It is whether one can prospectively
bound disagreement for a declared \emph{family of SNN execution semantics},
account for the gap between an emulator and physical hardware, and make the
result tight enough to direct repair.

We pursue four contributions:
\begin{enumerate}
    \item an operational, versioned execution-semantics object that fixes the
    digital choices omitted by a model graph;
    \item a hierarchy from exact mappings through sound per-input family analysis
    to simultaneous distribution and physical certificates;
    \item label-free certificate-directed calibration selected separately from
    the audit used for the final bound; and
    \item prospective validation on SpiNNaker2 and a new parameterized Virtex-7
    SRNN realization under a preregistered two-task, five-seed protocol.
\end{enumerate}

The intended claim is conditional and falsifiable. We will not claim that two
update equations are equivalent because a pilot trace happened to match. We will
not call an emulator-only result a physical certificate. We will not proceed with
the present framing if the static bound is mostly vacuous or physical conformance
dominates the total risk.


\section{Operational execution semantics}

Let a trained finite-horizon SNN be a graph $G_\theta$, an input event tensor
$x_{1:T}$, and a reference executor $F_{s_0}$. An execution semantics object $s$
fixes the timestep; forward or exponential Euler integration; whether threshold
evaluation occurs before or after integration; the complete integrate--threshold--
reset order; subtractive or reset-to-value behavior; state and weight precision,
scale, rounding, overflow, and saturation; synaptic and output delivery delays;
and deterministic or stochastic parameters.

For a current-based LIF unit with time constant $\tau$, the two supported
integration transitions are
\begin{align}
v'_{t+1} &= v_t + \Delta t(-v_t + I_t)/\tau,
    &&\text{forward Euler},\\
v'_{t+1} &= \alpha v_t + (1-\alpha)I_t,\quad
\alpha=e^{-\Delta t/\tau},
    &&\text{exponential Euler}.
\end{align}
Post-integration thresholding computes $z_t=\mathbf{1}[v'_{t+1}\geq\vartheta]$
and then either $v_{t+1}=v'_{t+1}-z_t\vartheta$ or
$v_{t+1}=z_t v_{\mathrm{reset}}+(1-z_t)v'_{t+1}$. Pre-integration thresholding
tests and resets $v_t$ before adding the current. These transitions are not
generally equal.

Fixed-point transitions are integer operations, not floating-point perturbations.
Each format names its signed bit width, fractional bits, rounding rule, and
overflow rule. Delays introduce explicit queue state. A one-step synaptic delay
therefore changes the finite horizon even if the eventual infinite event stream
would be a shift of the same sequence.

We write $\mathcal{S}$ for a bounded semantics family. Discrete axes are finite
sets. Continuous axes such as timestep or threshold scale are closed intervals.
Every semantics artifact and model is canonically serialized and hashed, so a
certificate cannot silently refer to a different executor.

\paragraph{Semantic validation.}
Every transition is checked on handcrafted traces and randomized scalar-versus-
vector differential tests. General equivalence is admitted only by symbolic
mapping. On a declared small finite input/state domain, exhaustive checking can
prove equality or return a minimal counterexample. This distinction separates a
global theorem, a finite-domain theorem, and an empirical observation.


\section{A four-level transport certificate}

\subsection{Exact and per-input certificates}

The first level searches a library of proved parameter transports. Structural
identity is the trivial member. Any nontrivial mapping must preserve the complete
transition, including reset and queues, and survive symbolic derivation,
exhaustive small-state checking, and adversarial trace search.

For a general bounded family, we unroll the recurrent network for $T$ steps and
propagate lower and upper state sets. Affine maps use sign-split interval
arithmetic. Fixed-point rounding is enclosed at integer bin boundaries; wrapping
that may cross the representable range expands to the full numeric domain.
Delays are queues of abstract states. At a threshold discontinuity, the analyzer
splits quiet and spiking branches, applies the appropriate reset, and conservatively
merges them. Discrete semantic variants are analyzed separately and their output
sets are joined.

For reference prediction $c_0$ and target logit intervals
$[\ell_c(x),u_c(x)]$, input $x$ is certified if
\begin{equation}
  \ell_{c_0}(x) > \max_{c\ne c_0} u_c(x).
\end{equation}
The condition proves an unchanged argmax for every semantics member represented
by the abstract family. Exact enumeration is also available when the declared
family itself is finite; a MILP checker is planned as a tighter small-network
oracle.

\subsection{Distribution certificate}

Static certification need not cover every input. On an untouched unlabeled audit
sample $x_1,\ldots,x_n$, define
\begin{equation}
  D_i(s)=\mathbf{1}[F_{s_0}(x_i)\ne F_s(x_i)].
\end{equation}
For $k_s=\sum_i D_i(s)$, we use the exact one-sided Clopper--Pearson limit
$U_s(\alpha)$. With $m$ target family members, Bonferroni correction at
$\alpha=\delta/m$ yields simultaneous coverage, and
$U_{\TRef,\Emu}=\max_s U_s(\delta/m)$ for a fixed but unknown declared target.

\begin{proposition}[Accuracy-change bound]
For arbitrary labels $Y$ and paired predictions from classifiers $f$ and $g$,
\[
|\Pr[f(X)=Y]-\Pr[g(X)=Y]|\leq \Pr[f(X)\ne g(X)].
\]
\end{proposition}
\begin{proof}
Pointwise, the two correctness indicators can differ only when the predictions
differ. Taking expectations and applying $|\mathbb{E}Z|\leq\mathbb{E}|Z|$
proves the statement.
\end{proof}

\subsection{Physical certificate}

An emulator cannot certify unmodeled hardware behavior. We therefore execute an
unlabeled canary and estimate paired emulator--hardware disagreement. Splitting
failure probability across the two estimates gives
\begin{equation}
U_{\mathrm{total}}=\min\{1,
U_{\TRef,\Emu}(\delta/2)+U_{\Emu,\Hw}(\delta/2)\}.
\end{equation}
Without the second term, the report is marked \emph{conditional on emulator}.
For nondeterministic hardware, the primary sample contains independent
$(x,\text{run})$ pairs. Repeated execution of selected inputs characterizes
run-to-run variance but cannot replace those pairs.

\section{Certificate-directed repair}

Repair follows a fixed order. First, a proved closed-form transport is applied
when one exists. Second, the primary method calibrates each neuron's threshold,
leak/time constant, bias, numeric scale, and permitted delay compensation without
labels. Third, an optional target-semantics fine-tune may use a declared few-shot
label budget.

Let $c_0(x)$ be the reference class and $m_s(x)$ its target logit margin over the
largest competing class. Rather than minimizing only mean logit distance, the
primary objective approximates uncertified prediction mass:
\begin{equation}
\mathcal{L}_{\mathrm{cert}}=
\frac{1}{|C|}\sum_{x\in C}w(x)\,\mathrm{softplus}(-m_s(x)/\gamma)
+\lambda_z d_z+\lambda_v d_v,
\end{equation}
where $C$ is the unlabeled calibration split, $w(x)$ emphasizes informative
reference margins, and $d_z,d_v$ compare spike and state traces. A vectorized
smooth-spike executor optimizes this objective with straight-through fixed-point
transitions; auditable hard-semantics coordinate search then refines and validates
the candidate.

No calibration sample enters the final audit. After a candidate is selected, its
model hash is frozen and a new certificate is built on disjoint IDs. The repair
report records permitted changes, data and label budgets, objective evaluations,
runtime, original and repaired hashes, and the post-repair certificate.

Comparisons include no repair, post-training quantization, global threshold
scaling, per-platform quantization-aware training, logit-only distillation, and
fully supervised target retraining. Calibration examples, label counts,
optimization evaluations, wall-clock/GPU cost, and hardware runs are matched or
reported explicitly.

The differentiable margin term is a proxy for eventual certification, not a
static certificate. We therefore evaluate separately whether a selected repair
increases sound-analysis coverage or reduces proof complexity. If it only
restores sampled decisions or labeled accuracy, we call it \emph{transport
repair}, not certificate restoration.

As a development-only proof-aware variant, we also penalize hidden units whose
hard target execution approaches a threshold guard. This local penalty is
selected on calibration data and evaluated by the independent sound analyzer;
it is not itself a verified multi-step margin. We stop this route before the
five-seed matrix if it fails to complete cells or outperform the unrepaired
model's proof depth.

\section{Physical realizations}

\paragraph{SpiNNaker2.}
The adapter pins the \texttt{py-spinnaker2} package and firmware, imports the same
NIR graph through the official \texttt{s2\_nir.from\_nir} route, and captures the
complete conversion configuration. Reset, discretization, and timing translations
are encoded in the execution-semantics sidecar rather than left as undocumented
adapter behavior. Each run stores spike outputs, predictions, sample IDs, pair
IDs, and repeated-run diagnostics.

\paragraph{Virtex-7.}
Our new digital target parameterizes state/weight widths, fractional scale,
integration, threshold order, reset rule, saturation or wrapping, and synaptic
and output delays. A bit-accurate integer oracle shares the RTL transition
contract. Verification compares every state and spike through randomized cocotb
tests, checks bounded reset/overflow properties, and finally compares captured
board traces. The bitstream, firmware, constraints, synthesis tools, and device
serial are immutable manifest fields.

For both devices, the reference, emulator, and physical executor receive
identical weights and event samples. A physical result is rejected if sample
pairing, semantics identity, firmware, or capture hashes cannot be reconstructed.


\section{Preregistered evaluation}

The primary event-native tasks are SHD with an SRNN and DVS Gesture with a
convolutional/recurrent SNN. N-MNIST with a feedforward SNN is a negative control
expected to be more transport-robust. SSC is deferred until all primary gates
pass. Five training seeds share frozen sample IDs and semantics variants.

The canonical training set is hash-partitioned into training, label-free repair
calibration, and unlabeled certificate audit splits. The canonical labeled test
set is sequestered. We vary reset, integration, threshold timing, precision,
rounding/saturation, and one-step delay separately, followed by four preregistered
high-risk combinations.

For every condition we will report actual absolute accuracy change, certificate
upper bound and slack, static certified-input fraction, sample complexity,
emulator--hardware disagreement, decisions at one-, two-, and five-point budgets,
repair recovery, remaining risk, label count, runtime, and hardware runs. Single-
axis experiments support causal attribution and provide counterexamples to
invalid equivalence claims.

\begin{table}[t]
\centering
\caption{Frozen primary matrix. Checkmarks denote required measurements, not
completed experiments.}
\begin{tabular}{lccc}
\toprule
Task & Reference--emulator & SpiNNaker2 & Virtex-7 \\
\midrule
SHD SRNN & \checkmark & \checkmark & \checkmark \\
DVS Gesture conv/SRNN & \checkmark & \checkmark & \checkmark \\
N-MNIST feedforward control & \checkmark & diagnostic & diagnostic \\
\bottomrule
\end{tabular}
\end{table}

The submission gate requires two physical backends, four conditions with more
than five points of uncorrected loss, no observed simultaneous-bound violation,
median slack at most three points, 90th-percentile slack at most eight points,
and at least 70\% label-free recovery across both tasks and backends. The method
must beat global threshold scaling and PTQ at matched cost. If more than 80\% of
inputs lack a nonvacuous static family certificate, or conformance dominates, we
pivot to a systems venue.

\TODO{Insert frozen five-seed results only after the prospective hardware matrix
and clean-machine replay pass.}


\section{Related work, assumptions, and failure cases}

NIR supplies interoperable primitives and explicitly documents that digital
implementations can diverge because of discretization, neuron models,
quantization, and nondeterminism \citep{pedersen2024nir}. Our work treats NIR as
the graph transport layer and adds a separate operational semantics and evidence
layer; it does not modify NIR's representation contract.

General machine learning has an extensive literature on disagreement and
label-free performance prediction. Disagreement-on-the-line studies when model
disagreement tracks performance under distribution shift \citep{lee2023dotl},
while disagreement discrepancy optimizes an auxiliary classifier to obtain error
bounds \citep{rosenfeld2023disagreement}. We cite these as prior art. Our basic
accuracy-change inequality is elementary. The proposed contribution is prospective
SNN execution-family certification, sound state analysis, physical conformance,
and repair under a separated audit protocol.

Formal SNN robustness work encodes spike dynamics in SMT and reports severe
scaling limits for rate-coded models; a more efficient direct search relies on
temporal encoding in which each neuron fires at most once
\citep{seong2024verification}. Our recurrent rate-coded setting permits repeated
spikes and varies execution parameters rather than inputs. Its explicit
branch-set failure is therefore consistent with, but not solved by, that line of
work.

The guarantee assumes a fixed finite horizon, immutable model and semantics,
paired audit inputs drawn independently from the deployment distribution, and
correct capture identities. Distribution drift after certification is outside
the claim. A finite enumerated family says nothing about an omitted backend
behavior. Interval merging may be sound but vacuous near repeated threshold
branches. Clopper--Pearson bounds can require large audit sets at one-point
budgets. A small emulator--hardware canary bound does not diagnose the cause of
conformance failures.

The scope is digital classification. We exclude analog mismatch, continual
learning, control, safety certification, and energy claims. A certificate bounds
absolute accuracy change, not class-conditional harm unless the sampling and
report are extended accordingly.


\bibliographystyle{plainnat}
\bibliography{references}

\appendix
\section{Artifact and proof obligations}

Every certificate stores the model, data, checkpoint, reference and target
semantics, code revision, split, seed, firmware, and bitstream hashes; sample
count; simultaneous confidence; static certified fraction; semantic,
conformance, and total bounds; budget verdict; and assumptions. Artifacts are
created once with exclusive filesystem semantics.

The software test matrix covers scalar/vector differential execution across both
integration and reset rules, multiple delays, and fixed-point formats; numeric
overflow; semantic hash stability; exact confidence limits; simultaneous
composition; small-domain counterexamples; split disjointness; repair
re-certification; CLI workflows; and physical manifest validation.

\section{Sample complexity at zero observed disagreement}

For zero disagreements in $n$ independent pairs, the one-sided
Clopper--Pearson upper limit is
\begin{equation}
U(\alpha)=1-\alpha^{1/n}.
\end{equation}
This expression makes the cost of small decision budgets explicit and is used to
plan canary size before hardware execution rather than after observing results.

\section{Proof-status vocabulary}

\begin{description}
\item[Structural identity.] Canonical semantics hashes and model transitions are
identical.
\item[Mapped equivalence.] A symbolic parameter map preserves every transition
under stated preconditions.
\item[Finite-domain proof.] Exhaustive states/inputs and horizon are recorded;
the conclusion does not extend beyond them.
\item[Static certificate.] A sound abstract set proves the argmax over one input
and the declared family.
\item[Statistical certificate.] Coverage depends on the paired audit sampling
assumptions.
\item[Physical certificate.] Statistical certificate plus a valid device canary
and immutable manifest.
\end{description}


\end{document}
